
the arrows displayed here are restriction homomorphisms, and one uses the fact that @@M0@@ is an isomorphism by the proper base change theorem (SGA 4 XII 5.5). I claim that \emph{the square} (7.13.2) \emph{is commutative}. The proof is essentially the same as that establishing the commutativity of (7.10.2); here it suffices to use the compatibility of the formation of Chern classes with inverse images for the morphisms @@M1@@ and @@M2@@.

\sgaStableTarget{sga6.e10.para.7-13-3}\textbf{7.13.3.} Note that the square (7.13.2) arises by passage to the projective limit from the analogous squares in which cohomology with coefficients in the @@M3@@-adic sheaf @@M4@@ is replaced by cohomology with coefficients in the torsion sheaves @@M5@@ (where @@M6@@ runs through the powers of @@M7@@). Thus the commutativity of (7.13.2) merely expresses the commutativity of the squares involving Chern classes mod. @@M8@@, for every power @@M9@@ of @@M10@@. Choose an algebraic closure @@M11@@ of @@M12@@, and hence an algebraic closure @@M13@@ of @@M14@@ (namely the residue field of the integral closure of @@M15@@ in @@M16@@). For every finite subextension @@M17@@ of @@M18@@, one can consider the analogous commutative diagram (for fixed @@M19@@) in the situation obtained from @@M20@@ by the base change @@M21@@, with @@M22@@ the normalization of @@M23@@ in @@M24@@. For this, note that, by Krull--Akizuki and since @@M25@@ is henselian, @@M26@@ is indeed a trait; moreover, one must suppose here that the fibers of @@M27@@ over @@M28@@ are geometrically regular, i.e., that @@M29@@ is \emph{smooth} over @@M30@@, so that this hypothesis is preserved by base change. As @@M31@@ varies, one thus obtains an inductive system of commutative squares. Passing to the inductive limit and using (SGA 4 VII 5.8) and (IV 3.2), one obtains a commutative square (cf. 7.16):
